\chapter{Conclusion}

\section{Conclusion}

This study aimed to explore the effectiveness of Vision Transformers (ViTs) for sleep stage classification, particularly focusing on improving model performance through various image transformation techniques, sampling frequencies, pre-trained models, and data augmentation methods. Through a series of experiments, we identified the optimal model and approach for classifying sleep stages using PSG signals. The ViT-B/16 model fine-tuned on the combined original and synthetic dataset with GASF and DWT at 128 Hz emerged as the best performer, demonstrating its ability to classify sleep stages with high accuracy and generalization capability. \\

Our findings align with previous studies highlighting the effectiveness of deep learning models, particularly ViTs, in handling complex time-series data like PSG signals. The use of synthetic data generated by DCGAN also proved beneficial, suggesting that GAN-based augmentation can help mitigate issues related to class imbalances and subject-to-subject variability. These results contribute to the growing body of research that emphasizes the potential of deep learning models, particularly ViTs, in enhancing the accuracy of sleep stage classification. \\

Despite the promising results, a few unexpected findings were observed, such as slight degradation in accuracy when using only synthetic data. While the GAN-augmented dataset led to a substantial improvement when combined with the original data, the performance of synthetic-only data was lower than expected. This highlights the need for further optimization in the GAN generation process to ensure more reliable synthetic data.

\section{Limitations}

While this study demonstrates the potential of Vision Transformers for sleep stage classification, several limitations should be acknowledged. One key limitation is the limited dataset scope—although the model was evaluated on DS1–DS8 datasets, its performance may not generalize to all sleep conditions or populations. Larger and more diverse datasets are needed to validate its robustness further. Another limitation is the quality of synthetic data generated by the GAN model. While synthetic data helped mitigate class imbalance, it was not entirely reliable when used independently, suggesting a need for more advanced GAN architectures to improve data fidelity. Additionally, the computational complexity of Vision Transformers presents a challenge for real-time deployment, particularly in wearable sleep monitoring devices, where efficiency and low power consumption are critical. \\

Furthermore, despite achieving high accuracy, the model's lack of interpretability remains a concern. Since deep learning models, including ViTs, function as black boxes, understanding their decision-making process is difficult. Future studies should incorporate explainability techniques such as attention visualization or SHAP analysis to enhance trust in clinical applications. Another limitation is the focus on EEG signals, as PSG data also includes other physiological signals like ECG, EOG, and EMG, which were not fully explored in this study. The integration of these additional modalities could potentially improve classification performance. Lastly, the study relied on fixed hyperparameters, without extensive optimization across different datasets. A more systematic hyperparameter tuning approach could further refine the model’s performance, ensuring better adaptation to varying sleep conditions. Addressing these limitations in future work could enhance the model’s reliability, interpretability, and practical applicability in real-world sleep monitoring scenarios.

\section{Further Work}

While the findings of this study provide valuable insights into the use of Vision Transformers for sleep stage classification, there are several areas where further research can build upon this work.

\begin{enumerate}
    \item \textbf{Improving GAN Generation Quality}: Future work could focus on improving the quality of synthetic data generated by DCGAN, potentially exploring advanced GAN architectures like StyleGAN or conditional GANs (cGANs) to reduce artifacts and blurring in the generated images.
    \item \textbf{Larger and More Diverse Datasets}: To further test the generalization capability of the model, future research could involve training and evaluating the model on a larger and more diverse set of datasets, beyond DS1–DS8, to assess its robustness across varied populations and sleep conditions.
    \item \textbf{Real-World Data Applications}: The practical implications of the model’s performance in real-world sleep studies should be explored. Evaluating the model on datasets from different clinical settings or using long-term sleep data from wearable devices would provide insights into the applicability of this approach in real-world environments.
    \item \textbf{Hybrid Models}: Investigating hybrid models combining ViTs with other deep learning techniques, such as Convolutional Neural Networks (CNNs) or Recurrent Neural Networks (RNNs), could potentially enhance performance, especially for sequential data.
    \item \textbf{Exploring Interpretability}: Although this study did not focus on interpretability, future research could look into explaining model predictions, which is particularly important in medical domains where understanding model decisions can be critical for clinical applications.
\end{enumerate}
By addressing these areas, future research can further advance the field of sleep stage classification, making it more accurate, reliable, and applicable for real-world use. \cite{mohomed2024nocturne}

 